\chapter{CƠ SỞ LÝ THUYẾT}
\section{Các khái niệm nền tảng}
\textit{Quyết định chọn trường} là mức độ ưu tiên và sẵn sàng đăng ký vào một cơ sở đào tạo sau khi cân nhắc các lựa chọn. \textit{Quyết định chọn ngành} là sự lựa chọn lĩnh vực đào tạo dựa trên nhận thức về bản thân, nội dung học và kết quả nghề nghiệp. Trong luận văn, biến phụ thuộc phản ánh ý định chắc chắn, ưu tiên nguyện vọng và sẵn sàng theo học; đây là đại diện gần của hành vi, không đồng nhất với nhập học thực tế.

\textit{Phù hợp cá nhân} thể hiện sự tương thích giữa sở thích, năng lực, giá trị và yêu cầu ngành. \textit{Triển vọng nghề nghiệp} gồm khả năng có việc, thu nhập, phát triển và nhu cầu nhân lực. \textit{Danh tiếng--chất lượng} là nhận thức về uy tín, chương trình, giảng viên và điều kiện đảm bảo chất lượng. \textit{Chi phí--hỗ trợ} bao gồm học phí, sinh hoạt phí, học bổng và khả năng chi trả. \textit{Ảnh hưởng xã hội} đến từ gia đình, bạn bè, giáo viên và người đã học. \textit{Truyền thông tuyển sinh} phản ánh tính đầy đủ, tin cậy và dễ tiếp cận của thông tin. \textit{Vị trí--điều kiện} bao gồm khoảng cách, an toàn, ký túc xá và môi trường sống.

\section{Các lý thuyết liên quan}
\subsection{Mô hình lựa chọn trường của Chapman}
Chapman \citep{chapman1981} chia tác nhân thành đặc điểm học sinh và ảnh hưởng bên ngoài. Đặc điểm học sinh gồm thành tích, kỳ vọng và nền tảng gia đình; ảnh hưởng bên ngoài gồm người quan trọng, đặc tính trường và hoạt động truyền thông. Khung này phù hợp để tổ chức biến nhưng cần mở rộng yếu tố phù hợp ngành và triển vọng nghề nghiệp.

\subsection{Mô hình ba giai đoạn}
Hossler và Gallagher \citep{hossler1987} mô tả lựa chọn đại học qua khuynh hướng, tìm kiếm và lựa chọn. Một thông điệp tuyển sinh có thể hiệu quả ở giai đoạn tìm kiếm nhưng ít ý nghĩa nếu học sinh chưa hình thành ý định học đại học. Do khảo sát cắt ngang học sinh lớp 12, nghiên cứu tập trung vào cuối giai đoạn tìm kiếm và giai đoạn lựa chọn.

\subsection{Lý thuyết hành vi có kế hoạch}
Theo Ajzen \citep{ajzen1991}, ý định được hình thành từ thái độ, chuẩn chủ quan và kiểm soát hành vi cảm nhận. Trong mô hình này, lợi ích nghề nghiệp và phù hợp cá nhân góp phần tạo thái độ; ảnh hưởng xã hội đại diện chuẩn chủ quan; tài chính và địa lý phản ánh một phần kiểm soát cảm nhận.

\subsection{Lý thuyết vốn con người và phù hợp người--môi trường}
Lý thuyết vốn con người xem giáo dục như đầu tư để tạo dòng lợi ích tương lai \citep{becker1993}, giải thích vai trò việc làm và chi phí. Tiếp cận phù hợp người--môi trường cho rằng kết quả tích cực xuất hiện khi đặc điểm cá nhân tương thích với môi trường học tập/nghề nghiệp \citep{holland1997}. Hai tiếp cận bổ sung nhau: học sinh vừa cân nhắc hiệu quả đầu tư, vừa tìm kiếm lựa chọn hợp với bản thân.

\section{Lược khảo nghiên cứu}
Nghiên cứu tại Việt Nam bằng phương pháp hỗn hợp cho thấy học sinh sử dụng các yếu tố khác nhau tùy khối ngành; ngày hội mở đặc biệt hữu ích với nhóm kỹ thuật, còn môi trường tiếng Anh có thể hấp dẫn nhóm mong muốn quốc tế hóa \citep{le2022}. Nghiên cứu tại An Giang dùng EFA và hồi quy để đánh giá quyết định chọn trường \citep{hai2023}. Các nghiên cứu gần đây tại ĐBSCL tiếp tục ghi nhận vai trò của cơ hội nghề nghiệp, danh tiếng, quan hệ xã hội và chi phí \citep{nguyen2023,nguyen2026}. Tổng hợp cho thấy không nên giả định một yếu tố thống trị ở mọi địa bàn; thang đo phải được hiệu chỉnh theo ngôn ngữ và điều kiện vùng.

\begin{table}[H]\centering\caption{Tổng hợp hướng nghiên cứu liên quan}\label{tab:review}
\begin{tabularx}{\textwidth}{p{2.7cm}p{3.1cm}X p{2.3cm}}\toprule
Nguồn & Bối cảnh/phương pháp & Nhóm yếu tố nổi bật & Khoảng trống\\\midrule
Chapman (1981) & Mô hình khái niệm & Cá nhân, người ảnh hưởng, đặc tính trường & Chưa nhấn mạnh ngành\\
Hossler và Gallagher (1987) & Mô hình quá trình & Khuynh hướng, tìm kiếm, lựa chọn & Không phải mô hình đo lường\\
Lê và cộng sự (2022) & Việt Nam, hỗn hợp & Nguồn thông tin, chương trình, khác biệt khối ngành & Chưa riêng ĐBSCL\\
Hải (2023) & An Giang, EFA/hồi quy & Đặc điểm trường, cơ hội, xã hội & Một tỉnh\\
Nguyễn và cộng sự (2023) & Học sinh từ ĐBSCL & Nhiều nhóm tác động đến chọn trường & Tập trung trường hơn ngành\\\bottomrule
\end{tabularx}\end{table}

\section{Mô hình và giả thuyết nghiên cứu}
\begin{figure}[H]\centering
\begin{tikzpicture}[font=\small,node distance=5mm and 25mm,box/.style={draw,rounded corners,minimum width=4.5cm,minimum height=8mm,align=center},result/.style={draw,rounded corners,fill=gray!12,minimum width=4.7cm,minimum height=16mm,align=center},arr/.style={-{Latex[length=2mm]},thick}]
\node[box] (fit) {Phù hợp cá nhân (FIT)};
\node[box,below=of fit] (car) {Triển vọng nghề nghiệp (CAR)};
\node[box,below=of car] (rep) {Danh tiếng--chất lượng (REP)};
\node[box,below=of rep] (cos) {Chi phí--hỗ trợ tài chính (COS)};
\node[box,below=of cos] (soc) {Ảnh hưởng xã hội (SOC)};
\node[box,below=of soc] (com) {Truyền thông tuyển sinh (COM)};
\node[box,below=of com] (loc) {Vị trí--điều kiện học tập (LOC)};
\node[result,right=35mm of cos] (dec) {Quyết định chọn\\trường/ngành (DEC)};
\foreach \x in {fit,car,rep,cos,soc,com,loc}\draw[arr] (\x.east) -- (dec.west);
\end{tikzpicture}
\caption{Mô hình nghiên cứu đề xuất}\label{fig:model}
\end{figure}

Các giả thuyết H1--H7 lần lượt phát biểu rằng FIT, CAR, REP, COS, SOC, COM và LOC tác động cùng chiều đến DEC. Riêng COS được đo theo hướng ``khả năng chi trả và mức hỗ trợ thuận lợi'', nên hệ số kỳ vọng dương.

\section{Thang đo nghiên cứu}
\begin{longtable}{p{1.2cm}p{9.6cm}p{2cm}}
\caption{Thang đo đề xuất}\label{tab:scale}\\\toprule Mã & Phát biểu rút gọn & Nguồn\\\midrule\endfirsthead
\toprule Mã & Phát biểu rút gọn & Nguồn\\\midrule\endhead
FIT1--4 & Phù hợp sở thích; phù hợp năng lực; phù hợp giá trị; tự tin theo học & Tổng hợp\\
CAR1--4 & Dễ có việc; thu nhập kỳ vọng; cơ hội phát triển; nhu cầu nhân lực & \citealp{becker1993}\\
REP1--4 & Uy tín; chương trình; giảng viên; kiểm định/cơ sở học tập & \citealp{chapman1981}\\
COS1--4 & Học phí phù hợp; sinh hoạt phí chấp nhận; học bổng; gia đình có khả năng & Tổng hợp\\
SOC1--4 & Cha mẹ; thầy cô; bạn bè; người đang/đã học & \citealp{ajzen1991}\\
COM1--4 & Thông tin đầy đủ; đáng tin; tư vấn hữu ích; kênh trực tuyến dễ dùng & \citealp{hossler1987}\\
LOC1--4 & Khoảng cách; giao thông; an toàn; ký túc xá/môi trường sống & Tổng hợp\\
DEC1--4 & Ưu tiên cao; sẵn sàng đăng ký; ít thay đổi; sẵn sàng theo học & Tổng hợp\\\bottomrule
\end{longtable}

\section{Tiêu chuẩn đánh giá thang đo}
Cronbach's Alpha từ 0,7 trở lên được xem là tốt trong nghiên cứu ứng dụng; tương quan biến--tổng nên lớn hơn 0,3 \citep{hair2019,nunnally1994}. EFA yêu cầu KMO $>0{,}5$, Bartlett có ý nghĩa, tải nhân tố ưu tiên $\geq0{,}5$, chênh lệch tải đủ rõ và tổng phương sai trích $\geq50\%$. Hồi quy được đánh giá bằng $R^2$ hiệu chỉnh, kiểm định F, hệ số chuẩn hóa, VIF, phần dư và các quan sát ảnh hưởng. Các ngưỡng là hướng dẫn, không thay thế phán đoán lý thuyết.

\section{Phân biệt chọn trường và chọn ngành}
Chọn ngành thường bắt đầu từ câu hỏi ``học gì và làm nghề gì'', trong khi chọn trường trả lời ``học ở đâu và trong môi trường nào''. Hai quyết định tương tác theo ba kiểu: ngành trước--trường sau; trường trước--ngành sau; và tối ưu đồng thời. Học sinh định hướng nghề rõ có xu hướng tìm trường đáp ứng ngành mong muốn. Ngược lại, học sinh ưu tiên thương hiệu, vị trí hoặc khả năng trúng tuyển có thể điều chỉnh ngành theo danh mục của trường. Bảng hỏi vì vậy yêu cầu người trả lời nghĩ đến một cặp trường--ngành cụ thể.

\section{Cơ chế tác động và biến kiểm soát}
FIT tác động thông qua hứng thú và tự hiệu quả; CAR thông qua lợi ích kỳ vọng; REP làm giảm bất định về chất lượng; COS và LOC quyết định tính khả thi; SOC cung cấp chuẩn mực và nguồn lực thông tin; COM giảm chi phí tìm kiếm. Các biến giới tính, học lực, thu nhập, khu vực, trình độ học vấn của cha mẹ và mức độ rõ ràng nghề nghiệp được cân nhắc làm biến kiểm soát.

\section{Ma trận giả thuyết và cơ chế}
\begin{longtable}{p{1.2cm}p{2.5cm}p{6.3cm}p{1.7cm}}\toprule Giả thuyết & Biến & Cơ chế kỳ vọng & Chiều\\\midrule\endfirsthead\toprule Giả thuyết & Biến & Cơ chế kỳ vọng & Chiều\\\midrule\endhead H1 & FIT & Tương thích làm tăng tự tin và cam kết & Dương\\H2 & CAR & Lợi ích nghề nghiệp làm tăng giá trị lựa chọn & Dương\\H3 & REP & Uy tín và chất lượng giảm rủi ro cảm nhận & Dương\\H4 & COS & Khả năng chi trả làm lựa chọn khả thi & Dương\\H5 & SOC & Sự ủng hộ làm tăng chuẩn chủ quan & Dương\\H6 & COM & Thông tin tốt làm giảm bất định & Dương\\H7 & LOC & Tiếp cận thuận lợi làm tăng kiểm soát cảm nhận & Dương\\\bottomrule\end{longtable}

\section{Nguy cơ sai lệch và biện pháp kiểm soát}
Sai lệch chọn mẫu được giảm bằng phân tầng địa bàn và loại trường; sai lệch không phản hồi được đánh giá bằng so sánh người trả lời sớm--muộn. Sai lệch phương pháp chung được hạn chế bằng ẩn danh, tách nhóm câu hỏi và ngôn ngữ trung tính. Kiểm định một nhân tố Harman chỉ được dùng như chẩn đoán sơ bộ, không phải bằng chứng duy nhất. Thiên lệch xã hội được giảm bằng việc khẳng định không có câu trả lời đúng sai.
